% GENERATED FILE. Edit canonical lessons, then run npm run generate:books.
% canonical-source-hash: fnv1a64:b381470df1e1a476
% canonical-lessons: RU-C350-posledstvie, RU-C350-sledstvie, RU-C350-otrazhatsya, RU-C350-effekt, RU-C350-v-rezultate

\chapter{Consequences}
\label{ch:ru-consequences}
\hlchaptermodality{350}

\begin{chapteropening}
\textbf{By the end of this chapter:} \emph{I can speak about consequences and results, say that something tells on something else, and describe an effect.}

Complete the last lesson of chapter 350: I can speak about consequences and results, say that something tells on something else, and describe an effect.
\end{chapteropening}

\section[\texorpdfstring{\ru{последствие} (poslédstviye)}{последствие (poslédstviye)}]{\ru{последствие} (poslédstviye) — a consequence}

\label{lesson:RU-C350-posledstvie}



\begin{quote}
Before the new one: say the Russian for influence, then the Russian for an impact.
\end{quote}

\subsection*{You'll want to know: \ru{последствие}}
\textbf{\ru{последствие}} (\emph{poslédstviye}) — \textquotedblleft{}a consequence\textquotedblright{}.

\begin{quote}
\ru{Это} \ru{решение} \ru{будет} \ru{иметь} \ru{серьёзные} \ru{последствия}.
\end{quote}

\begin{quote}
\emph{Éto reshéniye búdet imét seryóznyye poslédstviya.}
\end{quote}

\begin{quote}
This decision will have serious consequences.
\end{quote}

Neuter; mostly plural, \emph{\ru{последствия}}, and usually unpleasant.

\begin{cousinweb}
\textbf{\ru{последствие}} is built from \emph{\ru{после}} \textquotedblleft{}after\textquotedblright{} and \emph{\ru{след}} \textquotedblleft{}a track, a trail, a mark\textquotedblright{} (EtymDB), with the suffix \emph{-\ru{ствие}}: what comes after, in the wake of an event.
\end{cousinweb}

\begin{grammarlens}[title={What you've built}]
One more word for telling, explaining and describing at B1.
\end{grammarlens}

\subsection*{Your turn}
\begin{itemize}
\raggedright
  \item \emph{Say it:} \emph{poslédstviye}
  \item \emph{Say it:} \emph{poslédstviye}, once more
  \item \emph{Say it:} a sentence of your own with \emph{poslédstviye}
  \item \emph{Recall:} say the Russian for influence, then the Russian for an impact, then say \emph{poslédstviye} again
\end{itemize}

\begin{tcolorbox}[breakable,colback=teal!4,colframe=teal!35!black,title={Before you move on}]
What does \ru{последствие} mean? (\textquotedblleft{}A consequence\textquotedblright{}.) Say it once more.
\end{tcolorbox}
\section[\texorpdfstring{\ru{следствие} (slédstviye)}{следствие (slédstviye)}]{\ru{следствие} (slédstviye) — a result, an outcome; an investigation}

\label{lesson:RU-C350-sledstvie}



\begin{quote}
Before the new one: say the Russian for an impact, then the Russian for a consequence.
\end{quote}

\subsection*{You'll want to know: \ru{следствие}}
\textbf{\ru{следствие}} (\emph{slédstviye}) — \textquotedblleft{}a result, an outcome; an investigation\textquotedblright{}.

\begin{quote}
\ru{Его} \ru{болезнь} — \ru{следствие} \ru{стресса}.
\end{quote}

\begin{quote}
\emph{Yegó bolézn — slédstviye stréssa.}
\end{quote}

\begin{quote}
His illness is the result of stress.
\end{quote}

Neuter, stressed on the first syllable. In law, \emph{\ru{следствие}} is a criminal investigation.

\begin{cousinweb}
\textbf{\ru{следствие}} is built on \emph{\ru{след}} \textquotedblleft{}a track, a trail, a mark\textquotedblright{} (EtymDB) with the suffix \emph{-\ru{ствие}}, a Church Slavonic pattern: a result follows its cause like a track.
\end{cousinweb}

\begin{grammarlens}[title={What you've built}]
One more word for telling, explaining and describing at B1.
\end{grammarlens}

\subsection*{Your turn}
\begin{itemize}
\raggedright
  \item \emph{Say it:} \emph{slédstviye}
  \item \emph{Say it:} \emph{slédstviye}, once more
  \item \emph{Say it:} a sentence of your own with \emph{slédstviye}
  \item \emph{Recall:} say the Russian for an impact, then the Russian for a consequence, then say \emph{slédstviye} again
\end{itemize}

\begin{tcolorbox}[breakable,colback=teal!4,colframe=teal!35!black,title={Before you move on}]
What does \ru{следствие} mean? (\textquotedblleft{}A result, an outcome; an investigation\textquotedblright{}.) Say it once more.
\end{tcolorbox}
\section[\texorpdfstring{\ru{отражаться} (otrazhátsya)}{отражаться (otrazhátsya)}]{\ru{отражаться} (otrazhátsya) — to affect, to tell on; to be reflected}

\label{lesson:RU-C350-otrazhatsya}



\begin{quote}
Before the new one: say the Russian for a consequence, then the Russian for a result.
\end{quote}

\subsection*{You'll want to know: \ru{отражаться}}
\textbf{\ru{отражаться}} (\emph{otrazhátsya}) — \textquotedblleft{}to affect, to tell on; to be reflected\textquotedblright{}.

\begin{quote}
\ru{Усталость} \ru{отражается} \ru{на} \ru{работе}.
\end{quote}

\begin{quote}
\emph{Ustálost otrazháyetsya na rabóte.}
\end{quote}

\begin{quote}
Tiredness tells on your work.
\end{quote}

Imperfective; perfective \emph{\ru{отразиться}}. \emph{\ru{Отражаться} \ru{на}} + prepositional — to affect. Literally to be reflected: \emph{\ru{в} \ru{воде} \ru{отражается} \ru{луна}}.

\begin{cousinweb}
\textbf{\ru{отражаться}} is built from \emph{\ru{от}-} \textquotedblleft{}back\textquotedblright{} and the root of \emph{\ru{разить}} \textquotedblleft{}to strike\textquotedblright{} (EtymDB): light that strikes a surface and comes back is reflected, and so is a cause that comes back in its effects.
\end{cousinweb}

\begin{grammarlens}[title={What you've built}]
One more word for telling, explaining and describing at B1.
\end{grammarlens}

\subsection*{Your turn}
\begin{itemize}
\raggedright
  \item \emph{Say it:} \emph{otrazhátsya}
  \item \emph{Say it:} \emph{otrazhátsya}, once more
  \item \emph{Say it:} a sentence of your own with \emph{otrazhátsya}
  \item \emph{Recall:} say the Russian for a consequence, then the Russian for a result, then say \emph{otrazhátsya} again
\end{itemize}

\begin{tcolorbox}[breakable,colback=teal!4,colframe=teal!35!black,title={Before you move on}]
What does \ru{отражаться} mean? (\textquotedblleft{}To affect, to tell on; to be reflected\textquotedblright{}.) Say it once more.
\end{tcolorbox}
\section[\texorpdfstring{\ru{эффект} (effékt)}{эффект (effékt)}]{\ru{эффект} (effékt) — an effect}

\label{lesson:RU-C350-effekt}



\begin{quote}
Before the new one: say the Russian for a result, then the Russian for to affect.
\end{quote}

\subsection*{You'll want to know: \ru{эффект}}
\textbf{\ru{эффект}} (\emph{effékt}) — \textquotedblleft{}an effect\textquotedblright{}.

\begin{quote}
\ru{Новое} \ru{лекарство} \ru{дало} \ru{хороший} \ru{эффект}.
\end{quote}

\begin{quote}
\emph{Nóvoye lekárstvo dálo khoróshiy effékt.}
\end{quote}

\begin{quote}
The new medicine had a good effect.
\end{quote}

Masculine. \emph{\ru{Произвести} \ru{эффект}} — to make an impression.

\begin{cousinweb}
\textbf{\ru{эффект}} comes from Latin \emph{effectus} \textquotedblleft{}an effect\textquotedblright{} (EtymDB), from \emph{efficere} \textquotedblleft{}to bring about\textquotedblright{}, built from \emph{ex} \textquotedblleft{}out\textquotedblright{} and \emph{facere} \textquotedblleft{}to do, to make\textquotedblright{} (EtymDB; Lewis \& Short).
\end{cousinweb}

\begin{grammarlens}[title={What you've built}]
One more word for telling, explaining and describing at B1.
\end{grammarlens}

\subsection*{Your turn}
\begin{itemize}
\raggedright
  \item \emph{Say it:} \emph{effékt}
  \item \emph{Say it:} \emph{effékt}, once more
  \item \emph{Say it:} a sentence of your own with \emph{effékt}
  \item \emph{Recall:} say the Russian for a result, then the Russian for to affect, then say \emph{effékt} again
\end{itemize}

\begin{tcolorbox}[breakable,colback=teal!4,colframe=teal!35!black,title={Before you move on}]
What does \ru{эффект} mean? (\textquotedblleft{}An effect\textquotedblright{}.) Say it once more.
\end{tcolorbox}
\section[\texorpdfstring{\ru{в} \ru{результате} (v rezultáte)}{в результате (v rezultáte)}]{\ru{в} \ru{результате} (v rezultáte) — as a result}

\label{lesson:RU-C350-v-rezultate}



\begin{quote}
Before the new one: say the Russian for to affect, then the Russian for an effect.
\end{quote}

\subsection*{You'll want to know: \ru{в} \ru{результате}}
\textbf{\ru{в} \ru{результате}} (\emph{v rezultáte}) — \textquotedblleft{}as a result\textquotedblright{}.

\begin{quote}
\ru{В} \ru{результате} \ru{мы} \ru{опоздали} \ru{на} \ru{поезд}.
\end{quote}

\begin{quote}
\emph{V rezultáte my opozdáli na póyezd.}
\end{quote}

\begin{quote}
As a result, we missed the train.
\end{quote}

With a noun it takes the genitive: \emph{\ru{в} \ru{результате} \ru{аварии}}. On its own it opens a sentence.

\begin{cousinweb}
\textbf{\ru{в} \ru{результате}} is \emph{\ru{в}} \textquotedblleft{}in\textquotedblright{} with the prepositional of \emph{\ru{результат}} \textquotedblleft{}a result\textquotedblright{}, which is usually traced, through western languages, to Medieval Latin \emph{resultatum}, built on \emph{resultāre} \textquotedblleft{}to spring back\textquotedblright{} (Lewis \& Short).
\end{cousinweb}

\begin{grammarlens}[title={What you've built}]
One more word for telling, explaining and describing at B1.
\end{grammarlens}

\subsection*{Your turn}
\begin{itemize}
\raggedright
  \item \emph{Say it:} \emph{v rezultáte}
  \item \emph{Say it:} \emph{v rezultáte}, once more
  \item \emph{Say it:} a sentence of your own with \emph{v rezultáte}
  \item \emph{Recall:} say the Russian for to affect, then the Russian for an effect, then say \emph{v rezultáte} again
\end{itemize}

\begin{tcolorbox}[breakable,colback=teal!4,colframe=teal!35!black,title={Before you move on}]
What does \ru{в} \ru{результате} mean? (\textquotedblleft{}As a result\textquotedblright{}.) Say it once more.
\end{tcolorbox}
